% Part: first-order-logic
% Chapter: axiomatic-deduction
% Section: deduction-theorem-quantifiers

\documentclass[../../../include/open-logic-section]{subfiles}

\begin{document}

\olfileid{fol}{axd}{ddq}

\olsection{পরিমাণসূচক-সহ নিঃসরণ উপপাদ্য}

\begin{thm}[নিঃসরণ উপপাদ্য]
\ollabel{thm:deduction-thm-q} $\Gamma \cup \{!A\} \Proves !B$ হলে
$\Gamma \Proves !A \lif !B$।
\end{thm}

\begin{proof}
$\Gamma \cup \{!A\}$ থেকে~$!B$-এর !!{derivation}-এর দৈর্ঘ্যের ওপর
আবার আরোহ প্রয়োগ করি।

আরোহের ভিত্তির প্রমাণটি \olref[ded]{thm:deduction-thm}-এর প্রমাণের
ভিত্তির সঙ্গে অভিন্ন।

আরোহের ধাপে আবার ধরা যাক, $\Gamma \cup \{!A\}$ থেকে~$!B$-এর
!!{derivation} এমন ধাপ~$!B$ দিয়ে শেষ হয়, যা কোনো অনুমান-বিধি দিয়ে
সমর্থিত। অনুমান-বিধিটি মোডাস পোনেন্স হলে
\olref[ded]{thm:deduction-thm}-এর প্রমাণের মতোই এগোই। বিধিটি~\QR হলে
জানি, $!B \ident !C \lif \lforall[x][!D(x)]$ এবং $!C \lif !D(a)$
রূপের কোনো !!a{formula} !!{derivation}-এ আগে এসেছে, যেখানে $a$,
$!C$, $!A$ বা~$\Gamma$-তে আসে না। তাই আমাদের আছে
\begin{align*}
  \Gamma \cup \{!A\} & \Proves !C \lif !D(a),\\
  \intertext{এবং আরোহের অনুমান প্রযোজ্য; অর্থাৎ আমাদের আছে}
    \Gamma & \Proves !A \lif (!C \lif !D(a)).\\
  \intertext{আবার}
  & \Proves (!A \lif (!C \lif !D(a))) \lif ((!A \land !C) \lif !D(a))\\
  \intertext{এবং মোডাস পোনেন্স থেকে পাই}
  \Gamma & \Proves (!A \land !C) \lif !D(a).\\
  \intertext{আইগেনচল-শর্তটি এখনও প্রযোজ্য বলে এই !!{derivation}-এ \QR দিয়ে সমর্থিত একটি ধাপ যোগ করে পাই}
    \Gamma & \Proves (!A \land !C) \lif \lforall[x][!D(x)].\\
    \intertext{আবার আমাদের আছে}
    & \Proves ((!A \land !C) \lif \lforall[x][!D(x)]) \lif (!A \lif (!C \lif \lforall[x][!D(x)])),\\
    \intertext{তাই মোডাস পোনেন্স অনুসারে}
    \Gamma & \Proves !A \lif (!C \lif \lforall[x][!D(x)]),
\end{align*}
অর্থাৎ $\Gamma \Proves !A \lif !B$।
% BN-SRC-061 TARGET-CORRECTION-BEGIN
\par\smallskip\noindent\textnormal{\small\textbf{উৎস-সংশোধন (BN-SRC-061):} হিমায়িত ইংরেজি প্রদর্শনের দ্বিতীয় সর্বতঃপ্রমাণিত শর্তাধীন সূত্রে শেষের বাইরের ডান বন্ধনীটি অনুপস্থিত। বাংলা প্রদর্শনে $!A \lif (!C \lif \lforall[x][!D(x)])$ সম্পূর্ণ বন্ধ করা হয়েছে।} % BN-SRC-061 TARGET-CORRECTION-END
% BN-SRC-062 TARGET-CORRECTION-BEGIN
\par\smallskip\noindent\textnormal{\small\textbf{উৎস-সংশোধন (BN-SRC-062):} হিমায়িত ইংরেজি প্রমাণের শেষ বাক্যটি সদ্যপ্রাপ্ত $\Gamma \Proves !A \lif !B$-কে ভুল করে $\Gamma \Proves !B$ বলে। বাংলা পাঠে উপপাদ্যের প্রয়োজনীয় শর্তাধীন সিদ্ধান্তটি পুনরুদ্ধার করা হয়েছে।} % BN-SRC-062 TARGET-CORRECTION-END

$!B$ যদি~\QR দিয়ে সমর্থিত হয়, কিন্তু তার রূপ $\lexists[x][!D(x)]
\lif !C$ হয়, সেই ক্ষেত্রটি অনুশীলনী হিসেবে থাকল।
\end{proof}

\begin{prob}
  \olref[fol][axd][ddq]{thm:deduction-thm-q}-এর প্রমাণ সম্পূর্ণ করো।
\end{prob}

\end{document}
